

\begin{Exercice}[5 minutes]\textbf{Maxim}
Soit le code suivant. Quelle est sa complexité?

\begin{lstlisting}[style=pythonstyle]
def boucle(n):
    k = 0
    for i in range(n):
        for j in range(100):
            k += 1
    return k
\end{lstlisting}

\mcqlist{
 $O(\log n)$,
 $O(n)$,
 $O(n^2)$,
 $O(1)$
}

\begin{solution}
$O(n)$
\end{solution}

\end{Exercice}


\begin{Exercice}[5 minutes]\textbf{Maxim}
Soit le code suivant. Quelle est sa complexité?

\begin{lstlisting}[style=pythonstyle]
def boucle(n):
    k = 0
    for i in range(n):
        for j in range(i):
            k += 1
    return k
\end{lstlisting}

\mcqlist{
 $O(n^2)$,
 $O(n)$,
 $O(n\cdot \log n)$,
 $O(\log n)$
}

\begin{solution}
$O(n^2)$
\end{solution}

\end{Exercice}


\begin{Exercice}[5 minutes]\textbf{Maxim}
Soit le code suivant. Quelle est sa complexité?

\begin{lstlisting}[style=pythonstyle]
i = 1
while i < n:
    for j in range(n):
        print(i, j)
    i = i * 2
\end{lstlisting}

\mcqlist{
 $O(n^2)$,
 $O(n)$,
 $O(n\cdot \log n)$,
 $O(\log n)$
}

\begin{solution}
$O(n\cdot \log n)$
\end{solution}

\end{Exercice}

\begin{Exercice}[5 minutes]\textbf{Aurore}
Quelle est la complexité de la fonction suivante?

\begin{lstlisting}[style=javastyle]
public static void fonction(String word,int m) {
   int n = word.length()-1;
   while (n >= 0) {
       for (int i=0;i<m;i++) {
           System.out.println(word.charAt(n));
       }
       n--;
   }
}
\end{lstlisting}

\mcqlist{
    $O(n)$,
    $O(n^2)$,
    $O(n\cdot m)$,
    $O(n+m)$
}

\begin{solution}
$O(n\cdot m)$
\end{solution}

\end{Exercice}

\begin{Exercice}[5 minutes]\textbf{Aurore}
En applicant la fonction \textsc{mergeSort} telle que vue en cours sur la liste suivante, quelle est la paire de sous-listes sur laquelle merge() ne sera à aucun moment appelé [9,6,2,4,5,8]? 

\begin{lstlisting}[style=pythonstyle, caption={Merge sort — Python}]
def mergeSort(arr):
    if len(arr) <= 1:
        return arr
    
    mid = len(arr) // 2
    left = mergeSort(arr[:mid])
    right = mergeSort(arr[mid:])
    
    return merge(left, right)
\end{lstlisting}

\mcqlist{
    {$[6], [2]$},
    {$[2, 6, 9], [4, 5, 8]$},
    {$[4], [5, 8]$},
    {$[2], [4]$}
}

\begin{solution}
Option D: [2], [4]
\end{solution}
\end{Exercice}

% \begin{Exercice}[5 minutes]\textbf{Sean}

% Comment modifier tri-fusion pour que la liste soit triée par ordre décroissant:

% \mcqlist{
%     {Inverser la liste avant de commencer le tri.},
%     {Changer l'ordre des appels récursifs (traiter la partie droite avant la gauche)},
%     {Modifier le signe de comparaison dans l'étape de fusion.},
%     {Aucune de ces réponses}
% }

% \begin{solution}
% Option C.
% \end{solution}


% \end{Exercice}

\begin{Exercice}[5 minutes]\textbf{Sean}
Combien de fois la fonction \texttt{merge} va être appelée lors du tri-fusion sur une liste de $n$ éléments  

\mcqlist{
    $n$,
    $\log_2 n$,
    $n-1$,
    Cela dépend de la parité de la liste
}

\begin{solution}
Option C. Soit $f(n)$ la fonction qui donne le nombre d'appels à \texttt{merge}. On peut définir cette fonction récursivement suivant l'implémentation de \texttt{mergeSort}: $f(n) = 1 + 2\cdot f(\frac{n}{2})$. En remplaçant $f(n) = n-1$ on trouve que cette égalité est vraie. 
\end{solution}
\end{Exercice}

\begin{Exercice}[5 minutes]\textbf{Sean}
Quelle est la modification la plus efficace pour correctement supprimer les doublons de la liste pendant son tri avec l'algorithme de tri-fusion. 

\mcqlist{
    {Trier normalement, puis parcourir la liste finale pour retirer les éléments identiques.},
    {Vérifier si l'élément existe déjà dans la liste complète avant chaque insertion.},
    {Si les deux éléments comparés sont égaux dans \texttt{merge}, n'en ajouter qu'un seul au résultat et avancer les deux indices.},
    {Aucune de ces réponses}
}
\begin{solution}
Option C. 
\end{solution}
\end{Exercice}

\begin{Exercice}[5 minutes]\textbf{Léa}
Soient les affirmations suivantes:

\begin{itemize}
\item Si Alice prend un dessert, Bob en prend un aussi.
\item Chaque jour, soit Bob, soit Cathy, mais jamais les deux, prend un dessert.
\item Alice ou Cathy, ou les deux, prennent chaque jour un dessert.
\item Si Cathy prend un dessert, Alice fait de même.
\end{itemize}

Quelle combinaison de personnes prend un dessert si toutes les affirmations sont vraies?

\mcqlist{
    Alice seulement,
    Bob seulement,
    Cathy seulement,
    Alice et Bob uniquement
}

\begin{solution}
Option D: Alice et Bob uniquement. Cathy ne peut pas en prendre (car alors Alice et Bob devraient aussi, cela contredit soit Bob soit Cathy).
Donc Alice et Bob prennent un dessert.
\end{solution}

\end{Exercice}

\begin{Exercice}[5 minutes]\textbf{Léa}
Dans un ensemble E, considérons trois parties A, B, C. On sait que:

\begin{itemize}
\item $\{h, b\} \subset \overline{A} \cap \overline{B}$
\item $\{a, f\} \subset A \cup C$
\end{itemize}

Laquelle des propositions suivantes est vraie?

\mcqlist{
    {$h \in A$},
    {$b \in B$},
    {$a \in A \cup C$},
    {$f \notin C$}
}

\begin{solution}
$a \in A \cup C$.
\end{solution}

\end{Exercice}

\begin{Exercice}[5 minutes]\textbf{Léa}
Lesquelles des affirmations suivantes sont vraies concernant une classe abstraite Java.

\mcqlist{
    Peut être instanciée directement,
    Doit obligatoirement contenir au moins une méthode abstraite,
    Peut contenir des méthodes abstraites ou non,
    Ne peut pas contenir d’attributs
}

\begin{solution}
Option C. Une classe abstraite peut avoir zéro ou plusieurs méthodes abstraites, et des méthodes classiques.
\end{solution}

\end{Exercice}

\begin{Exercice}[5 minutes]\textbf{Léa}
Lesquelles des affirmations suivantes sont vraies concernant une interface Java.

\mcqlist{
    Les méthodes sont private par défaut,
    Les méthodes sont public abstract par défaut,
    Les variables sont privées par défaut,
    Les variables ne peuvent pas être finales
}

\begin{solution}
Option B. Les méthodes sont public abstract par défaut.
\end{solution}

\end{Exercice}

\begin{Exercice}[5 minutes]\textbf{Léa}
Quelle est la séquence correcte d’étapes du tri fusion?

\mcqlist{
        {
        Fusionner toute la liste \newline
        Diviser en deux moitiés \newline
        Trier séparément \newline
        Arrêter
    },
    {
        Diviser la liste en deux moitiés \newline
        Diviser encore chaque moitié jusqu’à obtenir des listes de taille 1 \newline
        Fusionner les petites listes pour former des listes triées \newline 
        Continuer les fusions jusqu’à obtenir une seule liste triée \newline
    },
    {
        Trouver le plus petit élément \newline
        Le placer au début \newline
        Répéter pour le second plus petit, etc. \newline
        Arrêter
    },
    {
        Trier d’abord la moitié droite \newline
        Trier ensuite la moitié gauche \newline
        Coller les deux moitiés sans les comparer \newline
        Arrêter
    }
}

\begin{solution}
Option B. 
\end{solution}

\end{Exercice}

\begin{Exercice}[5 minutes]\textbf{Léa}
Quelle est la complexité du code suivant supposant que la liste \texttt{L} contient $n$ éléments?

\begin{lstlisting}[style=pythonstyle]
def multiply(L):
    total = 1
    for x in L:
        total = total * x
    return total
\end{lstlisting}

\mcqlist{
    $O(1)$, $O(\log n)$, $O(n)$, $O(n^2)$
}


\begin{solution}
$O(n)$. 
\end{solution}

\end{Exercice}

\begin{Exercice}[5 minutes]\textbf{Maxim}

Quelle condition impérative une collection doit-elle respecter pour qu'on puisse y appliquer une recherche binaire (Binary Search)?

\mcqlist{
    Elle doit être triée., Elle doit être stockée dans une liste chaînée., Elle ne doit pas contenir de doublons., Elle doit contenir uniquement des entiers.
}

\begin{solution}
Option A. 
\end{solution}
\end{Exercice}

\begin{Exercice}[5 minutes]\textbf{Maxim}

Soit un arbre de recherche binaire vide dans lequel on insère la séquence suivante : {47,23,89,11,37,73,97}. Si on insère ensuite la valeur 31, où atterrit-elle?

\begin{center}
\includegraphics[width=0.6\textwidth]{img/tree.png}
\end{center}

\mcqlist{
    Comme enfant droit de 11., Comme enfant gauche de 73., Comme enfant gauche de 37., Comme enfant droit de 23.    
}

\begin{solution}
Option C. 
\end{solution}
\end{Exercice}

\begin{Exercice}[5 minutes]\textbf{Amina}
Soit la fonction suivante qui vérifie si une chaîne de caractères est un palindrome. Quelle est sa complexité temporelle supposant que la longueur de la chaîne de caractères est $n$?

\begin{lstlisting}[style=pythonstyle]
def est_palindrome(chaine: str) -> bool:
    chaine_nettoyee = "".join(filter(str.isalnum, chaine.lower()))
    n = len(chaine_nettoyee)
    
    gauche = 0
    droite = n - 1
    
    while gauche < droite:
        if chaine_nettoyee[gauche] != chaine_nettoyee[droite]:
            return False
            
        gauche += 1
        droite -= 1
        
    return True
\end{lstlisting}

\mcqlist{
    $O(n)$,  $O(n \log(n)$,  $O(1)$, $O(n^2)$
}

\begin{solution}
$O(n)$
\end{solution}

\end{Exercice}

\begin{Exercice}[5 minutes]\textbf{Amina}
Quelle notation asymptotique donne une borne supérieure asymptotique pour une fonction $f(n)$?


\mcqlist{
     $\Omega$-notation ,  $\Theta$-notation,  $O$-notation, $n_0$-notation
}

\begin{solution}
$O$-notation
\end{solution}

\end{Exercice}

\begin{Exercice}[5 minutes]\textbf{Sean}
Quelles affirmations sont correctes concernant les graphes denses?

\mcqlist{
    Un graphe complet est un graphe dense.,
    Un graphe connecté est un graphe dense.,
    Les listes d'adjacence sont plus efficaces pour représenter des graphes denses.,
    Les matrices d'adjacence sont plus efficaces pour représenter des graphes denses.
}

\begin{solution}
A, D.
\end{solution}
\end{Exercice}


\begin{Exercice}[5 minutes]\textbf{Sean}
Lesquelles des affirmations suivantes concernant l'algorithme de Kruskal sont vraies?

\mcqlist{
    L'algorithme de Kruskal utilise une approche gloutonne pour trouver l'arbre couvrant de poids minimal.,
    L'algorithme de Kruskal fonctionne aussi avec des poids négatifs.,
    L'algorithme de Krusal fonctionne sur des graphes non-connectés.,
     L'algorithme de Kruskal trouve l'unique solution optimale pour le problème de l'arbre couvrant de poids minimal.
}

\begin{solution}
B, E.
\end{solution}
\end{Exercice}

\begin{Exercice}[5 minutes]\textbf{Sean}
Lesquelles des affirmations suivantes concernant l'algorithme de parcours en largeur (BFS) sont vraies?

\mcqlist{
    L'algorithme BFS peut être utilisé pour trouver le plus court chemin dans un graphe non pondéré.,
    L'algorithme BFS utilise une pile (stack) pour garder la trace des nœuds à visiter.,
    {L'algorithme BFS visite tous les nœuds à une distance donnée avant de visiter les nœuds à la distance supérieure.},
    {Si il n'existe pas de chemin entre le nœud de départ et un autre nœud, l'algorithme BFS peut entrer dans une boucle infinie.},
    {Si le graphe contient des cycles, l'algorithme BFS ne peut pas être utilisé pour trouver le plus court chemin}.
}

\begin{solution}
A, C.
\end{solution}
\end{Exercice}

\begin{Exercice}[5 minutes]\textbf{Sean}
Combien d'appels à \textsc{UNION} sont nécessaires pour fusionner n ensembles disjoints en un seul ensemble?

\mcqlist{
    $n$,
    $n - 1$,
    $n + 1$,
    $2n$
}

\begin{solution}
Option B. L'arbre minimum couvrant a toujours $n-1$ arêtes pour $n$ sommets. Chaque appel à \textsc{UNION} crée une arête.  
\end{solution}
\end{Exercice}

%% TODO: Add Reda's questions here 

\begin{Exercice}[5 minutes]\textbf{Léa}
On applique une recherche séquentielle sur la liste \texttt{L = [5, 12, 7, 12, 20]} et on cherche 12.

Que renvoie l’algorithme de recherche séquentielle suivant?

\begin{lstlisting}[style=pythonstyle]
def recherche_sequentielle(liste, e):
    if e in liste:
        return liste.index(e)
    return -1 
\end{lstlisting}

\mcqlist{
1, 3, 12, (1, 3)
}

\begin{solution}
1
\end{solution}


\end{Exercice}


\begin{Exercice}[5 minutes]\textbf{Léa}
On cherche \texttt{key = 11} dans la liste triée \texttt{L = [2, 4, 6, 10, 13, 18]}. Que renvoie l'algorithme de recherche binaire suivant?

\begin{lstlisting}[style=pythonstyle]
def recherche_binaire(key, array):
    low = 0
    high = len(array) - 1

    while low <= high:
        mid = (low + high) // 2

        if key < array[mid]:
            high = mid - 1
        elif key > array[mid]:
            low = mid + 1
        else:
            return mid

    return -1
\end{lstlisting}

\mcqlist{
3, KeyError: 11, -1, 11
}

\begin{solution}
-1
\end{solution}

\end{Exercice}


\begin{Exercice}[5 minutes]\textbf{Léa}
On considère le graphe suivant, représenté par une liste d’adjacence:

\begin{lstlisting}
G = {
   'A': ['B','D'],
   'B': ['C'],
   'C': ['E'],
   'D': ['E'],
   'E': []
}
\end{lstlisting}

L'algorithme de BFS: 

\begin{lstlisting}[style=pythonstyle]
def bfs(graph, start):
    visited = list(start) # liste des sommets visités
    queue = [start] # liste des sommets à visiter
    order = [] # pour stocker l'ordre de sommets
    u = queue.pop(0) 
    order.append(u)
    neighbors = graph[u] # liste des voisins de u
    for v in neighbors : 
        if v not in visited : 
            queue.append(v)
            visited.append(v)
return order 
\end{lstlisting}

Que retourne \texttt{BFS(G, 'A')}?

\mcqlist{
{$[$'A','B','D','C','E'$]$ }, {$[$'A','D','B','C','E'$]$}, {$[$'A','B','C','D','E'$]$}, {$[$'A','D','E','B','C'$]$}
}

\begin{solution}
Option A. 
\end{solution}

\end{Exercice}

\begin{Exercice}[5 minutes]\textbf{Léa}
Soit le graphe décrit par la liste d'adjacence (adjacency list) ci-dessous:

\begin{lstlisting}
G = {
  "A": [("B", 3), ("C", 1)],
  "B": [("A", 3), ("C", 2), ("D", 4)],
  "C": [("A", 1), ("B", 2), ("D", 5)],
  "D": [("B", 4), ("C", 5)]
}
\end{lstlisting}

Quel ensemble d’arêtes forme l’arbre couvrant minimal obtenu avec Kruskal?

\mcqlist{
{('A','C',1), ('B','C',2), ('B','D',4)},
 {('A','C',1), ('A','B',3), ('B','D',4)},
 {('A','C',1), ('B','C',2), ('A','B',3)},
{('A','B',3), ('B','C',2), ('C','D',4)}
}

\begin{solution}
Option A. 
\end{solution}

\end{Exercice}



\begin{Exercice}[5 minutes]\textbf{Sean}
Soit l'algorithme de construction d'un arbre KD (KD-tree) avec $K=3$, et un ensemble de points en 3D sous la forme $(x, y, z)$.
Si le premier axe de division est l'axe $X$, le deuxième $Y$ et le troisième $Z$, quel axe sera utilisé lors du troisième appel à \textsc{BuildKDTree} étant donné qu'il y a 10 points distincts?

\begin{algorithm}[H]
\caption{BuildKDTree($\text{points}, \text{depth}, k$)}
\label{alg:kdtree-build}
\begin{algorithmic}[1]
\If{$\text{points} = \emptyset$}
    \State \textbf{return} $\text{NIL}$
\EndIf
\State $\text{axis} \leftarrow \text{depth} \ \% \ k$ \Comment{Select axis based on current depth}
\State $\text{points} \leftarrow \text{SORT}(\text{points}, \text{key} = \text{axis})$ \Comment{Sort by current axis}
\State $\text{median\_idx} \leftarrow \text{INTEGER}(\text{length}(\text{points}) / 2 )$ \Comment{Index of median list item}
\State $\text{median} \leftarrow \text{KDNode}(\text{points}[\text{median\_idx}]$)
\State $\text{left} \leftarrow \text{points}[1 \ldots \text{median\_idx} - 1]$ \Comment{Left half of the points list}
\State $\text{right} \leftarrow \text{points}[\text{median\_idx} + 1 \ldots \text{length}(\text{points})]$ \Comment{Right half of the points list}
\State $\text{median.left} \leftarrow \text{BuildKDTree}(\text{left}, \text{depth} + 1, k)$
\State $\text{median.right} \leftarrow \text{BuildKDTree}(\text{right}, \text{depth} + 1, k)$
\State \textbf{return} median
\end{algorithmic}
\end{algorithm}

\mcqlist{
    X, Y, Z
}

\begin{solution}
Z
\end{solution}

\end{Exercice}

\begin{Exercice}[5 minutes]\textbf{Sean}
Lesquelles des affirmations suivantes concernant les quadtrees sont vraies?

\mcqlist{
    {Les Quadtrees sont toujours des arbres équilibrés.},
    {Chaque nœud interne d'un Quadtree a exactement quatre enfants.},
    {Les Quadtrees sont principalement utilisés pour représenter des données en trois dimensions.},
    {Avec une profondeur de $n$ (hauteur de l'arbre), un Quadtree peut représenter jusqu'à $2^n$ régions distinctes.},
    {Avec une profondeur de $n$ (hauteur de l'arbre), un Quadtree peut représenter jusqu'à $4^n$ régions distinctes. }
}

\begin{solution}
B, E. 
\end{solution}
 

\end{Exercice}



\begin{Exercice}[5 minutes]\textbf{Aurore}
Laquelle des affirmations suivantes est fausse concernant les algorithmes de recherche?

\mcqlist{
    {La complexité d'un algorithme de recherche binaire est de O(log(n)).},
    {La performance de la recherche binaire est meilleure que celle de la recherche séquentielle pour de grandes collections.},
    {La recherche par arbre de recherche binaire est une recherche dichotomique.},
    {La recherche séquentielle requiert une collection triée.},
    {Toutes les affirmations sont correctes. }
}
\begin{solution}
Option D. 
\end{solution}


\end{Exercice}
